% CAPÍTULO 5 — CONSIDERAÇÕES FINAIS

\chapter{Considerações Finais}
\label{cap:consideracoes}

Este trabalho se propôs a responder se é possível construir uma plataforma analítica reprodutível sobre os microdados da RAIS e do CAGED, com baixo custo operacional, e o que esses dados dizem sobre o mercado formal de tecnologia no Brasil. As duas faces da pergunta foram respondidas, e a resposta à segunda depende da primeira: foi a construção cuidadosa da camada de dados que permitiu enxergar coisas que agregados publicados não mostram.

\section{Contribuições}

A \textbf{contribuição técnica} é um \textit{pipeline} completo, do servidor público ao \textit{dashboard}, que trata 56,4~GB de microdados heterogêneos sem servidor próprio nem custo de infraestrutura, é executado automaticamente toda semana e entrega consulta interativa sobre a série inteira. Quatro decisões sustentam esse resultado: o formato colunar particionado, que permite descartar dados sem lê-los; o motor analítico \textit{in-process}, que dispensa servidor e importação prévia; a consolidação de arquivos por partição, que reduziu 2.111 objetos a 628 e levou a contagem completa dos vínculos da RAIS a 0,4~segundo; e a execução em nuvem, com o repositório público de dados como o próprio \textit{data lake} e etapas de integração contínua encadeadas por dependência real, que reprocessam só os anos com dado novo e podem ser repetidas uma a uma quando falham.

A \textbf{contribuição metodológica} está na forma de delimitar o setor. Em lugar de escolher entre definir tecnologia pela atividade da empresa ou pela ocupação do trabalhador, o trabalho preserva as duas lentes como rótulos independentes em cada registro. Essa escolha é o que torna possível o achado central: a maior parte dos profissionais de TI está fora das empresas de TI, e existe um contingente expressivo de pessoas em empresas de TI exercendo outras ocupações, com remuneração muito inferior. Nenhum recorte de lente única enxerga os dois grupos.

A \textbf{contribuição analítica} são as respostas às cinco perguntas, sintetizadas no \autoref{quadro:sintese}. Duas merecem destaque pelo que contrariam. A primeira é que crescer não melhorou a remuneração relativa: o emprego mais que dobrou e a mediana em salários mínimos ficou onde estava. A segunda é a distinção estrutural entre os eixos de desigualdade --- o hiato de gênero é inteiramente de retorno, com parcela explicada negativa, enquanto o racial é majoritariamente de acesso ---, distinção que implica instrumentos de política diferentes.

Há ainda uma \textbf{contribuição sobre a qualidade da fonte}, que não estava prevista. A auditoria encontrou defeitos concretos nos dados publicados pelo MTE e um erro de método próprio, relatados na \autoref{sec:achados-qualidade}. Nenhum deles produz mensagem de erro, e todos produzem resposta errada --- o que os torna especialmente úteis de documentar.

\section{Limitações}

As limitações de fonte foram declaradas na \autoref{sec:limitacoes-metodo} e merecem ser retomadas. A RAIS e o CAGED cobrem apenas o mercado formal, o que exclui a contratação como pessoa jurídica, comum no setor. A localização é a do estabelecimento declarante, e não a do trabalho --- distorção que o trabalho remoto agrava e que o dado público não permite corrigir. A CBO é atualizada lentamente, e o próprio trabalho encontrou código em uso que nenhum dicionário oficial documenta.

Três limitações de método também precisam ficar registradas. A primeira é que a parcela não explicada da decomposição salarial não mede discriminação: ela é um limite superior para o que a discriminação poderia explicar. A segunda é que a tábua de vida de período descreve o risco vigente em um ano, e não a trajetória de uma coorte --- a leitura de permanência precisa dessa ressalva. A terceira é que o Prophet, previsto na proposta, não foi avaliado: a família SARIMA foi adotada e validada, mas a comparação entre as duas abordagens ficou de fora.

\section{Estado Atual do Trabalho e Pendências}
\label{sec:pendencias}

A migração para a nuvem resolveu duas das pendências registradas na fase local: o bloqueio das bibliotecas de modelagem pela política de integridade de código do Windows, porque os modelos passaram a ser calculados em executores Linux, e a falta de espaço em disco, porque o \textit{data lake} deixou de ser local. Ela também mudou a natureza de uma terceira: não existe mais uma etapa de republicação a cumprir à mão, porque a reconstrução do site com os dados novos é a última etapa do próprio \textit{pipeline}. O \autoref{quadro:pendencias} relaciona o que permanece.

\begin{quadro}[htbp]
    \centering
    \caption{Pendências para a versão final}
    \label{quadro:pendencias}
    \footnotesize
    \renewcommand{\arraystretch}{1.3}
    \begin{tabularx}{\textwidth}{%
        >{\raggedright\arraybackslash\hsize=0.62\hsize}X
        >{\raggedright\arraybackslash\hsize=1.38\hsize}X}
        \toprule
        \rowcolor{black!8}
        \textbf{Pendência} & \textbf{Situação} \\
        \midrule
        Levar a correção de 2022 para a nuvem &
        A correção da ingestão parcial de 2022 (\autoref{sec:achados-qualidade}) foi aplicada à camada tratada local, e os agregados da RAIS e do mapa usados nas tabelas deste texto foram reconstruídos sobre ela. A Silver da RAIS publicada, porém, é anterior à correção --- e é ela que alimenta a Gold e o \textit{dashboard}. Como a verificação incremental compara a presença de anos, e não o conteúdo, um ano corrigido não é detectado sozinho: 2022 precisa ser reprocessado sob demanda, pela opção de reprocessamento forçado que a etapa da Silver passou a oferecer, depois da qual a Gold é reconstruída. A diferença está medida: regeradas a partir do repositório, oito das doze tabelas deste texto saem idênticas, e as duas da RAIS que diferem (\autoref{tab:estoque-serie} e \autoref{tab:linhas-camada}) diferem só em 2022, pelos 446.267 vínculos ausentes; essas duas permanecem com os valores corrigidos. No \textit{dashboard}, 2022 aparece sinalizado como subestimado. Afetados até o reprocessamento: os agregados da RAIS desse ano, o \textit{nowcast} e a linha de 2022 da série do hiato. \textbf{Não} afetados: a previsão do saldo mensal, a análise de sobrevivência (calculada sobre 2024), os perfis de município (janela de 2020 a 2025) e a análise de sensibilidade sobre São Paulo, que compara 2007 com 2025. \\
        Meses novos do CAGED na Silver &
        A verificação incremental da Silver do CAGED comparava a presença de \textbf{anos}: depois que o primeiro mês de um ano era traduzido, os seguintes nunca eram, porque o ano já constava. Por isso julho e agosto de 2026 estão na Bronze desde 8 de outubro e ainda não na Silver. A verificação foi corrigida para comparar também os meses, e a próxima execução precisa traduzi-los. \\
        Execução completa da Gold em nuvem &
        A consolidação das tabelas do CAGED e a unificação das duas gerações, etapa própria do \textit{pipeline}, foram executadas com sucesso em nuvem. Na execução agendada seguinte, as três etapas da Gold que leem a camada tratada falharam depois de 2 a 11 minutos de processamento, e a causa ainda precisa ser identificada no registro da execução. Até uma execução completa, a Gold publicada é a da última execução bem-sucedida. \\
        Substituição de partições na Silver em nuvem &
        Nas primeiras execuções em nuvem da Silver do CAGED, os arquivos de cada partição foram gravados fora da pasta da tabela e com nome diferente do usado pela camada migrada --- 505 arquivos na raiz do repositório. Eles não chegaram a ser lidos, porque as consultas leem só a pasta de cada tabela, mas também não atualizaram a camada; e, gravados na pasta certa, teriam contado a partição em dobro (\autoref{sec:publicacao}). A publicação da Silver, nas duas bases, foi refeita para \textbf{substituir} o ano reconstruído: cada partição vira um arquivo com o nome da camada consolidada, e tudo o que existia daquele ano --- inclusive as sobras na raiz --- sai no mesmo \textit{commit}. A substituição do ano inteiro é segura porque a partição é o ano e o mês do arquivo de origem, e não a competência: reconstruir um ano lê todos os arquivos dele. Se algum arquivo de origem falha, nada é publicado. Falta a primeira execução com a correção. \\
        Verificação de reprodutibilidade da RAIS &
        O teste de recontagem a partir do dado bruto (\autoref{sec:reprodutibilidade}) foi feito para o CAGED. O equivalente para a RAIS ainda não foi executado. \\
        Confronto com os painéis oficiais &
        A comparação dos agregados com os painéis do PDET, prevista na proposta, foi feita de forma pontual e precisa ser sistematizada por unidade federativa e ano. \\
        \bottomrule
    \end{tabularx}
    \\[4pt]
    Fonte: Autoria própria.
\end{quadro}

\section{Trabalhos Futuros}

Três extensões decorrem diretamente do que foi construído.

A primeira é medir \textbf{retenção por gênero} cruzando as duas bases: o CAGED informa quem entra, a RAIS informa quem permanece. A comparação entre a composição das admissões e a do estoque permitiria distinguir um problema de entrada de um problema de permanência --- pergunta que nenhuma das duas bases responde isolada.

A segunda é a \textbf{decomposição regional} do crescimento, no estilo \textit{shift-share}, separando quanto do crescimento de cada estado se explica pela composição setorial e quanto por dinamismo próprio. O dado já está na camada; falta o método.

A terceira é de engenharia: documentar a linhagem do \textit{pipeline} de forma navegável. A migração resolveu metade do problema --- o grafo de dependências entre as etapas é exibido pela plataforma de integração contínua a cada execução ---, e a procedência já existe no dado, linha a linha. Falta a mesma visão entre as tabelas dentro da Gold. Ferramentas de transformação declarativa geram esse grafo automaticamente, e a camada Gold --- composta de agregados em SQL --- é a parte do projeto que se beneficiaria dessa adoção, sem necessidade de reescrever a ingestão ou a camada tratada.

Por fim, a plataforma foi construída para ser reaproveitada. O desenho --- ingestão sem transformação, tradução por dicionários oficiais com procedência, recorte declarado e auditado, camada analítica em formato aberto, execução automatizada e dados tratados públicos --- não tem nada de específico do mercado de tecnologia. Aplicá-lo a outro recorte ocupacional, como engenharia, saúde ou educação, ou a outra base governamental brasileira de grande volume, é questão de trocar a definição do recorte e os dicionários de tradução. O primeiro passo nessa direção já foi dado: uma tabela de referência com a hierarquia oficial dos códigos de CBO e de CNAE, e uma coluna de categoria preenchida para o recorte de TI, foi publicada junto dos demais conjuntos de dados. Hoje o recorte ainda é lido de listas no código; quando o \textit{pipeline} passar a lê-lo dessa tabela, um recorte novo poderá ser definido marcando códigos nela, sem alterar as etapas de processamento.
